\apendice{Conjunto de Perguntas da Avaliação Experimental}
\label{ap:perguntas}

O Quadro~\ref{quad:perguntas-avaliadas} apresenta as 21 perguntas do conjunto de teste avaliadas no Capítulo~\ref{cap:resultados}, indicando o identificador, o documento de origem, a categoria da consulta, a página física de referência no PDF curado, o tipo de componente esperado no gabarito e o enunciado. Em seguida, o Quadro~\ref{quad:perguntas-excluidas} relaciona as três perguntas de Estatística que integravam o conjunto original de 24 questões de teste, mas cujas páginas de referência ficaram fora do recorte de 500 páginas processadas.

\begin{quadro}[htbp]
\centering
\caption{As 21 perguntas de teste avaliadas nas quatro configurações (E1 a E4).}
\label{quad:perguntas-avaliadas}
\begingroup
\scriptsize
\setlength{\tabcolsep}{3.5pt}
\renewcommand{\arraystretch}{0.96}
\begin{tabularx}{\textwidth}{|l|l|l|c|l|X|}
\hline
\textbf{ID} & \textbf{Doc.} & \textbf{Cat.} & \textbf{Pág.} & \textbf{Comp.} & \textbf{Pergunta} \\
\hline
\texttt{calculus\_v2\_q002} & Cálculo & Texto & 63 & \texttt{plain\_text} & Como o livro classifica a descontinuidade de $f(x)=(x^2-4)/(x-2)$ em $x=2$, e por quê? \\
\hline
\texttt{calculus\_v2\_q003} & Cálculo & Visual & 3 & \texttt{figure} & Qual formato geral e qual ponto mínimo aparecem no gráfico de $k(x)=x^2$ da Figura 2.3? \\
\hline
\texttt{calculus\_v2\_q004} & Cálculo & Visual & 13 & \texttt{figure} & Na Figura 2.12, qual função apresenta assíntota vertical em $x=2$ e valores positivos crescendo dos dois lados? \\
\hline
\texttt{calculus\_v2\_q005} & Cálculo & Visual & 101 & \texttt{figure} & Na Figura 3.8, qual reta representa a velocidade instantânea em um instante: a secante ou a tangente? \\
\hline
\texttt{calculus\_v2\_q008} & Cálculo & Fórmula & 110 & \texttt{isolate\_formula} & Como é definida $f'(x)$ pelo limite da razão incremental na Equação 3.9? \\
\hline
\texttt{calculus\_v2\_q009} & Cálculo & Multimodal & 7 & \texttt{figure} & Na aproximação de uma área sob uma curva mostrada nas Figuras 2.8 e 2.9, o que ocorre quando os retângulos se tornam mais estreitos? \\
\hline
\texttt{graphs\_v2\_q002} & Grafos & Texto & 41 & \texttt{plain\_text} & O que representa uma matriz de adjacência de um grafo e qual é sua dimensão? \\
\hline
\texttt{graphs\_v2\_q003} & Grafos & Texto & 82 & \texttt{plain\_text} & De acordo com o lema do aperto de mão, qual relação existe entre a soma dos graus dos vértices e a quantidade de arestas? \\
\hline
\texttt{graphs\_v2\_q005} & Grafos & Visual & 3 & \texttt{figure} & Na Figura 1.3, quais vértices são unidos pela aresta $e_5$? \\
\hline
\texttt{graphs\_v2\_q007} & Grafos & Visual & 94 & \texttt{figure} & Entre os grafos (a) e (b) da Figura 5.10, qual não satisfaz a condição suficiente de Ore, embora ainda possua ciclo hamiltoniano? \\
\hline
\texttt{graphs\_v2\_q008} & Grafos & Multimodal & 4 & \texttt{figure} & Na Figura 1.4, por que a aresta $e_6$ tem apenas $v_3$ como extremidade na tabela de incidência? \\
\hline
\texttt{graphs\_v2\_q010} & Grafos & Multimodal & 41 & \texttt{figure} & Por que a matriz de adjacência apresentada na Figura 2.5 é simétrica em relação à diagonal principal? \\
\hline
\texttt{physics\_v2\_q002} & Física & Visual & 8 & \texttt{figure} & Na Figura 3.7, em qual intervalo de tempo a posição do objeto permanece constante? \\
\hline
\texttt{physics\_v2\_q003} & Física & Visual & 81 & \texttt{figure} & Na Figura 4.23, para onde aponta a aceleração centrípeta em relação à trajetória circular? \\
\hline
\texttt{physics\_v2\_q006} & Física & Tabela & 77 & \texttt{table} & Segundo a Tabela 4.1, qual é a aceleração centrípeta de um satélite em órbita geossíncrona? \\
\hline
\texttt{physics\_v2\_q008} & Física & Fórmula & 77 & \texttt{isolate\_formula} & Qual é a relação entre a frequência angular $\omega$ e o período $T$ do movimento circular uniforme? \\
\hline
\texttt{physics\_v2\_q009} & Física & Multimodal & 8 & \texttt{figure} & Com base nas Figuras 3.7 e 3.8, qual é a velocidade do objeto entre 1,0 s e 2,0 s e como isso aparece no gráfico da posição? \\
\hline
\texttt{physics\_v2\_q010} & Física & Multimodal & 107 & \texttt{figure} & Na Figura 5.10, por que a situação em que o guincho puxa o automóvel produz aceleração maior do que a situação em que dois estudantes o empurram? \\
\hline
\texttt{statistics\_v2\_q006} & Estatística & Tabela & 8 & \texttt{table} & Na Tabela 3 de pesos dos potinhos de canela, qual intervalo de classe possui a maior frequência absoluta? \\
\hline
\texttt{statistics\_v2\_q009} & Estatística & Fórmula & 6 & \texttt{isolate\_formula} & Qual expressão é usada para calcular a amplitude total $A$ de um conjunto de dados? \\
\hline
\texttt{statistics\_v2\_q010} & Estatística & Multimodal & 9 & \texttt{figure} & A partir do histograma da Figura 10 e da Tabela 3, qual classe de peso concentra mais potinhos e qual é sua frequência relativa? \\
\hline
\end{tabularx}
\endgroup
\legend{Fonte: Elaborado pelo autor.}
\end{quadro}

\begin{quadro}[htbp]
\centering
\caption{As 3 perguntas de teste de Estatística excluídas antes da avaliação por dependerem de páginas fora do recorte de 500 páginas.}
\label{quad:perguntas-excluidas}
\begingroup
\scriptsize
\renewcommand{\arraystretch}{1.15}
\begin{tabularx}{\textwidth}{|l|l|c|l|X|}
\hline
\textbf{ID} & \textbf{Categoria} & \textbf{Pág.} & \textbf{Comp.} & \textbf{Pergunta} \\
\hline
\texttt{statistics\_v2\_q002} & Texto & 63 & \texttt{plain\_text} & Qual é a diferença entre erro do tipo I e erro do tipo II em um teste de hipóteses? \\
\hline
\texttt{statistics\_v2\_q003} & Evid. visual & 101 & \texttt{figure} & Que relação geral entre volume de água e volume de polpa é sugerida pelo diagrama de dispersão da Figura 34? \\
\hline
\texttt{statistics\_v2\_q007} & Tabelas & 101 & \texttt{table} & Segundo a Tabela 9 dos cocos verdes, qual é o teor de cálcio do fruto 3? \\
\hline
\end{tabularx}
\endgroup
\legend{Fonte: Elaborado pelo autor.}
\end{quadro}
